\scanpage{351}
% @ head of sheet
\mspage{75}
% @ top, lines 1–2
$\underline{\mathrm{Fl}}\,\mathrm{Cat}$. Ceci posé, l'axiome $\mathrm{Loc}\,9$ est le suivant

% @ upper third, lines 3–4
\[
\begin{aligned}
&\underline{\mathrm{Loc}}\ 9:\quad W\text{ est une partie \uline{accessible} de }\mathrm{Fl}(\mathrm{Cat})\\
&=\mathrm{Ob}\,\underline{\mathrm{Fl}}(\mathrm{Cat}).
\end{aligned}
\]
% @ middle, lines 5–11
5.3. Rappelons que cette notion un peu technique se simplifie considérablement, si on suppose que $W$ satisfait l'axiome des limites $\mathrm{Loc}(8)$. Alors $W$ est accessible \unsure{sss|ssi} il existe une \uline{petite} partie $W_{\unsure{\bullet|0|\circ}}$ de $\mathrm{Fl}(\mathrm{Cat})$, telle que $W$ soit la ``$L_1$-enveloppe'' de $W_0$, i.e.
% @ lower middle, lines 12–16
la plus petite partie de $\mathrm{Fl}(\mathrm{Cat})$ contenant $W_\bullet$ et \inserted{$L_1$-}stable \deleted{par} (i.e stable par $\underrightarrow{\lim}$ filtrantes). \deleted{\illegible{1}} S'il en est ainsi, \inserted{il est clair que} on peut \deleted{\unsure{trouver un toujours}} \inserted{prendre $W_\bullet$ de} la forme
% @ lower third, cardinal formula and right-hand insertion
\[
W_{\unsure{\pi|\kappa}}=W\cap\mathrm{Fl}_{\pi}(\mathrm{Cat})
\quad \inserted{\begin{aligned}&\text{où }\unsure{\kappa|\pi}\text{ est un}\\&\text{cardinal infini}\\&\text{convenable}\end{aligned}}
\]
% @ foot, last four lines
où $\mathrm{Fl}_{\kappa}(\mathrm{Cat})\overset{\mathrm{def}}{=}\mathrm{Fl}(\underline{\mathrm{Cat}}_{\pi})$, $\underline{\mathrm{Cat}}_{\pi}$ désignant la sous-catégorie \inserted{\unsure{str.}} pleine de $\mathrm{Cat}$ formée des \deleted{catégories} objets $X$ tels que $\operatorname{card} X\leq\pi$ (et par définition, $\operatorname{card} X=\operatorname{card}\mathrm{Fl}\,X$).
% archive sheet number: 350
